\input{../article_base.tex}
\setcounter{secnumdepth}{2}

\begin{document}

\section{הרצאה --- משפט ז'ורדן־ברוור ובורסוק־אולם}

\subsection{רקע}
נזכיר מה זה $I_2(f, X)$ ולמה זה מוגדר היטב.
\begin{corollary}
	אם $f : X \to Y$ העתקה חלקה ו־$\partial f : \partial X \to Y$ טרנסוורסלית ל־$Z \subseteq Y$ אז קיימת העתקה $g : X \to Y$ הומוטופית ל־$f$ כך ש־$\partial f = \partial g$ וגם $g \pitchfork Z$.
\end{corollary}
בניסוח אחר נקבל שאם $h : \partial X \to Y$ היא העתקה טרנסוורסלית ל־$Z \subseteq Y$ ו־$h$ ניתנת להרחבה לאיזושהי העתקה $X \to Y$, אז בפרט היא יכולה להתרחב להעתקה טרנסוורסלית ל־$Z$ בכל $X$.
\begin{definition}[אינדקס חיתוך מודולו 2]
	נניח ש־$X$ יריעה קומפקטית ו־$f : X \to Y$ העתקה חלקה טרנסוורסלית ל־$Z \subseteq Y$ סגורה, ונניח גם ש־$\dim X + \dim Z = \dim Y$.
	במקרה זה $f^{-1}(Z)$ היא תת־יריעה סגורה ממימד $0$ של $X$, כלומר היא קבוצה סופית.
	נגדיר את אינדקס החיתוך מודולו 2 של $f$ ו־$Z$, $I_2(f, Z)$, להיות מספר הנקודות ב־$f^{-1}(Z)$ מודולו 2. \\
	עבור העתקה חלקה כללית $g : X \to Y$ נבחר העתקה חלקה $f$ הומוטופית ל־$g$ וטרנסוורסלית ל־$Z$ ונגדיר $I_2(g, Z) = I_2(f, Z)$.
\end{definition}
נקבל ש־$I_2(g, Z)$ מוגדרת היטב בשל המשפט הבא,
\begin{theorem}
	אם $f_0, f_1 : X \to Y$ הומוטופיות וטרנסוורסליות ל־$Z$, אז $I_2(g, Z) = I_2(f, Z)$.
\end{theorem}
אם $X \subseteq Y$ אנו מקבלים את המקרה המיוחד שבו גם $i : X \hookrightarrow Y$ מקבלת אינדקס חיתוך מודולו 2, והוא כמובן משומר, נגדיר זאת.
\begin{definition}[אינדקס חיתוך מודולו 2 של תתי־יריעות]
	אם $X, Z \subseteq Y$ שתי תת־יריעות ממימד משלים, כלומר $\dim X + \dim Z = \dim Y$, וכן $X$ קומפקטית, אז נגדיר $I_2(X, Z) = I_2(i, Z)$.
\end{definition}
הדוגמה מעמוד 93.

עוד מקרה מיוחד שנרצה לדבר עליו הוא המקרה שבו $\dim X = \dim Y$ ו־$f : X \to Y$ חלקה.
אם נניח ש־$Y$ קשירה נקבל ש־$I_2(X, \{ y \})$ שווה לכמעט כל $y \in Y$.
\begin{definition}[דרגה מודולו 2]
	תהי $X$ קומפקטית ו־$Y$ קשירה, $f : X \to Y$ חלקה.
	נגדיר את הדרגה מודולו 2 של $f : X \to Y$ להיות הערך שווה לכמעט כל $I_2(f, \{ y \})$.
	נסמן את הדרגה ב־$\deg_2 f$.
\end{definition}
דוגמה של מעגל יחידה והפונקציה $z \mapsto z^n$. \\
ניצוק תוכן להגדרה זו עם המשפט הבא,
\begin{theorem}[שימור דרגה מודולו 2 על־ידי הומוטופיה]
	אם $f, g : X \to Y$ העתקות חלקות והומוטופיות, $X$ קומפקטית ו־$Y$ קשירה,
	אז,
	\[
		\deg_2 f = \deg_2 g
	\]
\end{theorem}
\begin{definition}[על־משטח]
	יריעה $X \subseteq \RR^n$ תיקרא על־משטח (Hypersurface) אם $\dim X = n - 1$.
\end{definition}
תהי $X \subseteq \RR^n$ על־משטח ו־$f : X \to \RR^n$ חלקה.
אנו מעוניינים להבין בהינתן נקודה $z \in Y$ לשאול כמה $f$ מתלפפת סביב $z$.
בהתאם נגדיר,
\[
	u(x)
	= \frac{f(x) - z}{|f(x) - z|}
.\]
ונקבל פונקציה אשר לכל $x$ מחזירה לנו את הכיוון כווקטור ב־$S^{n - 1}$ מ־$f(x)$ ל־$z$.
אנו טוענים שהליפוף שקול למספר הפעמים ש־$u$ מחזירה כיוון מסוים.
ממשפט החיתוך (שצריך לצטט פה קודם) אנו מסיקים ש־$u : X \to S^{n - 1}$ מצביעה פעמיים מודולו 2 כמעט תמיד אותה כמות של פעמים, ולמעשה הגדרנו מספר זה להיות בדיוק $\deg_2 u$.
\begin{definition}[אינדקס ליפוף מודולו 2]
	נגדיר את אינדקס הליפוף מודולו 2 ונסמן $W_2(f, z) = \deg_2 u$ להיות דרגת פונקציית הכיוון של $f$ מ־$z$.
\end{definition}
כמה דוגמות

\subsection{ז'ורדן־ברוור}
עתה נגדיר את המשפט שיעסיק אותנו בשעה הקרובה.
\begin{theorem}[ז'ורדן־ברוור]
	יהי $X \subseteq \RR^n$ על־משטח קשיר וקומפקטי,
	אז $X^C = \RR^n \setminus X = D_0 \cup D_1$ עבור שתי קבוצות פתוחות וזרות.
	$D_0$ החלק ''החיצוני'' למשטח ו־$D_1$ החלק ''הפנימי'' למשטח.
	בנוסף $\bar{D}_1$ היא יריעה קומפקטית המקיימת $\partial \bar{D}_1 = X$.
\end{theorem}
\begin{proof}
	TODO
\end{proof}

\subsection{ובורסוק־אולם}
\begin{theorem}[ובורסוק־אולם]
	תהי $f : S^k \to \RR^{k + 1}$ העתקה חלקה כך שהיא לא שולחת אף נקודה לראשית,
	ונניח ש־$f$ מקיימת את תנאי הסימטריה,
	\[
		\forall x \in S^k,\ 
		f(-x) = - f(x)
	.\]
	אז $W_2(f, 0) = 1$.
\end{theorem}
\begin{proof}
	TODO
\end{proof}

\end{document}
